% sections/coverpage.tex
% Cover page for BYU Mentored Research Grant Proposal

\thispagestyle{empty}
\begingroup
\setlength{\parskip}{0.2em}
\titlespacing*{\subsection}{0pt}{0.6em}{0.2em}

\begin{center}
{\Large \textbf{Mentored Research Grant (MRG)}}
\end{center}

% --- Project Title ---
\subsection*{Project Title}
Bayesian Optimization of Multi-Material 3D-Printed Tensegrity Structures for Energy Absorption

% --- Principal Investigator(s) ---
\subsection*{Principal Investigator(s)}
\textbf{PI Name:} Jeffrey R.\ Hill \\
\textbf{PI Department:} Mechanical Engineering \\
\textbf{Co-PI:} Sterling G.\ Baird, Department of Mechanical Engineering

% --- Abstract ---
\subsection*{Abstract}

We propose a two-year mentored research program in which 2 undergraduate
students design and optimize multi-material 3D-printed tensegrity-inspired
structures for maximizing energy absorption while limiting permanent deformation.  Tensegrity-inspired
structures---assemblies of rigid compression members (PLA) and flexible tension
members (TPU)---offer tunable mechanical
responses and high strength-to-weight ratios.  Using a multifidelity Bayesian
optimization framework, students will combine finite element simulations
(low-fidelity) with physical experiments (high-fidelity) to efficiently navigate
the high-dimensional design space.  Students will gain hands-on experience in
CAD design, multi-material 3D printing, impact testing, computational modeling,
and data-driven optimization.  Expected outcomes include undergraduate
conference presentations (UCUR, ASME IDETC), a peer-reviewed journal
submission, and development of marketable technical skills.  This project
provides a rich mentored learning environment at the intersection of structural
mechanics, advanced manufacturing, and machine learning. Undergraduates will lead both the computational and experimental components of the project, progressing from structured training in Year 1 to independent research leadership and peer mentoring in Year 2.

% --- Number of Students ---
\subsection*{Number of Students to be Mentored}
\textbf{Number of Undergraduate Students:} 2 \\
\textbf{Number of Graduate Students:} 1 (co-mentor, in support of creating the mentoring environment)

% --- Budget ---
\subsection*{Budget}

\vspace{-0.5em}
\begin{center}
\begin{tabular}{@{}lrrr@{}}
\toprule
\textbf{Item} & \textbf{Year 1} & \textbf{Year 2} & \textbf{Total} \\
\midrule
Undergraduate wages    & \$9,000  & \$9,000  & \$18,000 \\
Graduate wages         & \$1,250  & \$1,250  & \$2,500  \\
Supplies               & \$2,250  & \$750    & \$3,000  \\
Travel                 & \$0      & \$1,500  & \$1,500  \\
\midrule
\textbf{Total}         & \textbf{\$12,500} & \textbf{\$12,500} & \textbf{\$25,000} \\
\bottomrule
\end{tabular}
\end{center}
\vspace{-0.5em}

% --- Relationship to External Funding ---
\subsection*{Relationship to External Funding}

The PIs do not currently hold external funding for tensegrity or architected
materials research.  This MRG provides the primary support for establishing a
mentored undergraduate research program in this area.  If successful, the
preliminary data may inform a future NSF proposal; however, the MRG is
self-contained and its outcomes do not depend on external funding.

\endgroup
\newpage
